\documentclass[11pt]{article}

\usepackage[margin=1in]{geometry}
\usepackage{amsmath}
\usepackage{enumitem}
\usepackage{xcolor}
\usepackage{hyperref}
\usepackage{listings}

\hypersetup{
  colorlinks=true,
  linkcolor=blue,
  urlcolor=blue
}

\lstdefinestyle{code}{
  basicstyle=\ttfamily\small,
  breaklines=true,
  frame=single,
  columns=fullflexible
}
\lstset{style=code}

\title{Atmosphere, Surface Water, and Hydrothermal Environment Architecture}
\author{LIFEORIG working notes}
\date{\today}

\begin{document}
\maketitle

\section{Purpose}

This note summarizes the current architecture discussion for connecting
planetary parameters, stellar forcing, atmospheric disequilibrium chemistry,
surface water availability, and possible hydrothermal vent environments.

The main design principle is:

\begin{quote}
The physical solver should use physical parameters and computed states. World
labels such as Earth-like, super-Earth, mini-Neptune, ocean world, or Titan-like
should be assigned only after the solver produces its result.
\end{quote}

\section{Relevant Code Paths}

\begin{description}[leftmargin=3.7cm]
  \item[Input parser] \texttt{src/input\_data/read\_input.py}
  \item[Chemistry input/result data] \texttt{src/chemical\_env/chem\_env\_data.py}
  \item[Chemistry dispatcher] \texttt{src/chemical\_env/chemical\_environment.py}
  \item[Atmosphere solver driver] \texttt{src/atmosphere\_solver/atm\_struct\_driver.py}
  \item[Planet data] \texttt{src/planet\_params/planetary\_params.py}
  \item[Stellar data] \texttt{src/stellar\_params/stellar\_data.py}
  \item[Black-body radiation] \texttt{src/stellar\_params/black\_body\_radiation.py}
  \item[Stellar spectrum plot] \texttt{src/utilities/plot\_stellar\_spectrum.py}
  \item[Volcanic-rock water model] \texttt{src/environment/volcanic\_rocks/water\_level.py}
  \item[pH estimator] \texttt{src/environment/pH\_estimate.py}
  \item[Atmosphere test runner] \texttt{TEST\_SCRIPTS/run\_atmdyn\_solver.sh}
  \item[Titan test runner] \texttt{TEST\_SCRIPTS/run\_titan.sh}
\end{description}

\section{Input Files and Modes}

\subsection{Atmospheric disequilibrium input}

The script

\begin{lstlisting}
TEST_SCRIPTS/run_atmdyn_solver.sh
\end{lstlisting}

writes a full \texttt{input.json} into:

\begin{lstlisting}
TESTS/ATMOSPHERE/input.json
\end{lstlisting}

The relevant chemistry mode is:

\begin{lstlisting}
"exo_chemistry": {
    "mode": "layered_disequilibrium",
    "atomic_abundances": {
        "H": 1.0,
        "He": 0.1,
        "C": 1e-4,
        "O": 5e-3,
        "N": 1e-5
    }
}
\end{lstlisting}

The stellar input currently contains:

\begin{lstlisting}
"stellar_data": {
    "spectral_class": "G2V",
    "star_temperature": { "value": 5778, "units": "K" },
    "star_radius": { "value": 6.957e8, "units": "m" }
}
\end{lstlisting}

The stellar radius is given in meters because Pint did not recognize
\texttt{solar\_radius} in the input. The value \texttt{6.957e8 m} is the nominal
solar radius.

\section{Current Architecture}

\subsection{Input parsing}

\texttt{read\_input.py} parses:

\begin{itemize}
  \item planetary parameters into \texttt{PlanetaryEnvironmentParams};
  \item stellar parameters into \texttt{StellarParams};
  \item local environment parameters into \texttt{local\_env\_data};
  \item chemistry input into \texttt{ChemEnvInput}.
\end{itemize}

The summaries should belong to their data objects. Therefore
\texttt{PlanetaryEnvironmentParams} and \texttt{StellarParams} expose summary
methods, and \texttt{read\_input.py} only calls those methods.

\subsection{Chemistry dispatcher}

\texttt{SimulateChemEnv} dispatches chemistry modes:

\begin{lstlisting}
if mode == "local_equilibrium":
    run EasyChem
if mode == "layered_disequilibrium":
    run AtmosphSolver
\end{lstlisting}

The desired separation is:

\begin{itemize}
  \item \texttt{ChemEnvInput} remains chemistry-oriented.
  \item Stellar data is needed by the atmospheric solver, not by local
        equilibrium chemistry.
  \item \texttt{AtmosphSolver} consumes the stellar object when computing
        radiation.
\end{itemize}

\section{Stellar Black-Body Radiation}

The black-body spectral radiance is implemented in:

\begin{lstlisting}
src/stellar_params/black_body_radiation.py
\end{lstlisting}

The function is:

\begin{lstlisting}
blackbody_B_lambda(wavelength, temperature)
\end{lstlisting}

It returns Planck spectral radiance:

\[
B_{\lambda}(T) =
\frac{2hc^2}{\lambda^5}
\frac{1}{\exp\left(\frac{hc}{\lambda k_B T}\right)-1}
\]

with units:

\[
\mathrm{W \, m^{-3} \, sr^{-1}}.
\]

\subsection{Stellar flux at the planet}

In \texttt{AtmosphSolver}, the stellar intensity is evaluated over a wavelength
grid:

\begin{lstlisting}
self.wavelength_grid = Q_(np.linspace(100.0, 1000.0, 1000), "nanometer")
self.stellar_B_lambda = self.stellar_data.B_lambda(self.wavelength_grid)
\end{lstlisting}

The range \(100-1000\) nm covers UV, visible, and near-infrared. It is a good
first range for a Sun-like star, but it does not include the full infrared.

The base spectral flux at orbital distance is:

\[
F_{0,\lambda}
= \Omega_\star B_\lambda(T_\star)
\left(\frac{R_\star}{a}\right)^2,
\]

where the solid angle factor cancels the steradian in \(B_\lambda\). In code,
this is represented by multiplying by a solid-angle quantity:

\begin{lstlisting}
stellar_disk_solid_angle = Q_(np.pi, "steradian")
self.F0 = (
    stellar_disk_solid_angle
    * self.stellar_B_lambda
    * (self.stellar_data.radius / self.planet_data.orbital_distance) ** 2
).to("W / AU**2 / m")
\end{lstlisting}

The intended unit for \texttt{F0} is:

\[
\mathrm{W \, AU^{-2} \, m^{-1}},
\]

because it has been integrated over solid angle.

\section{Atmospheric Solver Output}

The atmosphere solver should eventually produce a structured atmospheric result,
including:

\begin{itemize}
  \item wavelength grid and stellar forcing;
  \item atmospheric layers;
  \item pressure and temperature profile;
  \item chemistry per atmospheric layer;
  \item surface temperature and pressure;
  \item surface partial pressures, especially \(p_{\mathrm{H_2O}}\);
  \item surface water phase and water availability.
\end{itemize}

The solver should be independent of labels such as ``Earth-like'' or ``ocean
world''. Those labels are diagnostics assigned after the physical state is
computed.

\section{World Classification}

After the atmosphere dynamics is solved, world type can be assigned from final
computed quantities:

\begin{itemize}
  \item planet mass and radius;
  \item bulk density or inferred volatile envelope;
  \item surface temperature and pressure;
  \item atmospheric composition;
  \item water partial pressure;
  \item liquid/ice/vapor water state;
  \item estimated surface water coverage.
\end{itemize}

The classifier should not control the solver. Instead:

\begin{lstlisting}
solver output -> classifier -> label
\end{lstlisting}

Example labels:

\begin{itemize}
  \item dry rocky world;
  \item wet rocky world;
  \item Earth-like;
  \item super-Earth;
  \item water world;
  \item ice world;
  \item steam world;
  \item mini-Neptune;
  \item Titan-like methane world.
\end{itemize}

\section{Surface Water and Volcanic Rocks}

The volcanic-rock water model is in:

\begin{lstlisting}
src/environment/volcanic_rocks/water_level.py
\end{lstlisting}

The intended split is:

\begin{itemize}
  \item the atmosphere solver computes surface atmospheric/surface water state;
  \item \texttt{water\_level.py} maps that state into local volcanic-rock water
        availability.
\end{itemize}

At present, \texttt{water\_level.py} estimates a water factor from atmospheric
\(\mathrm{H_2O}\) partial pressure:

\begin{lstlisting}
derive_atmospheric_h2o_water_factor(...)
\end{lstlisting}

This is a downstream consumer. It should not solve the atmosphere.

\section{pH Estimation}

The pH estimator lives in:

\begin{lstlisting}
src/environment/pH_estimate.py
\end{lstlisting}

Its purpose is to estimate water pH from chemistry if liquid water is available,
including volcanic-rock settings and later hydrothermal settings.

The current first-order logic is:

\begin{itemize}
  \item consume partial pressures from a \texttt{ChemEnvResult};
  \item dissolve acid/base gases with Henry-law approximations;
  \item estimate acid/base contributions to \([\mathrm{H^+}]\) or
        \([\mathrm{OH^-}]\);
  \item return a \texttt{pHResult}.
\end{itemize}

Current species include:

\begin{itemize}
  \item \(\mathrm{CO_2}\), acid;
  \item \(\mathrm{NH_3}\), base;
  \item \(\mathrm{H_2S}\), acid.
\end{itemize}

\section{Aqueous/Ocean Chemistry Backend}

For ocean or hydrothermal chemistry, the suitable backend is not EasyChem.
EasyChem is for gas-phase atmospheric equilibrium.

The practical aqueous/geochemical backend should be PHREEQC, likely through
\texttt{phreeqpy}, or later Reaktoro if a heavier backend is needed.

The backend input should include:

\begin{itemize}
  \item temperature;
  \item pressure if needed;
  \item pH guess;
  \item redox state, \(pe\) or \(Eh\);
  \item dissolved totals or species concentrations;
  \item gas partial pressures for gas-water exchange;
  \item minerals and phases for water-rock interaction.
\end{itemize}

The planned coupling is iterative:

\begin{enumerate}
  \item atmosphere solver computes surface chemistry;
  \item pH estimator gives an initial pH guess;
  \item aqueous backend computes consistent aqueous speciation and pH;
  \item updated aqueous chemistry feeds back into the environment model;
  \item repeat until pH and composition converge.
\end{enumerate}

\section{Hydrothermal Vent Model}

Hydrothermal vents should be modeled as a gradient between two reservoirs.

\subsection{Ambient ocean reservoir}

Inputs come from atmosphere/surface/ocean coupling:

\begin{itemize}
  \item \(T_{\mathrm{ocean}}\);
  \item \(P_{\mathrm{ocean}}\);
  \item \(pH_{\mathrm{ocean}}\);
  \item dissolved species;
  \item redox state;
  \item gas exchange with the atmosphere.
\end{itemize}

\subsection{Vent reservoir}

Inputs come from geophysics and water-rock interaction assumptions:

\begin{itemize}
  \item \(T_{\mathrm{vent}}\);
  \item \(P_{\mathrm{vent}}\);
  \item \(pH_{\mathrm{vent}}\);
  \item reduced species such as \(\mathrm{H_2}\), \(\mathrm{CH_4}\),
        \(\mathrm{H_2S}\);
  \item metal and mineral chemistry such as Fe, Mg, Si, sulfides;
  \item rock/mineral phases;
  \item flow rate and water-rock ratio.
\end{itemize}

\subsection{Gradients}

The vent can be approximated as a one-dimensional reactive channel:

\[
T(x) = T_{\mathrm{ocean}} +
\left(T_{\mathrm{vent}} - T_{\mathrm{ocean}}\right)
\exp(-x/L_T),
\]

\[
pH(x) = pH_{\mathrm{ocean}} +
\left(pH_{\mathrm{vent}} - pH_{\mathrm{ocean}}\right)
\exp(-x/L_{pH}),
\]

\[
C_i(x) = C_{i,\mathrm{ocean}} +
\left(C_{i,\mathrm{vent}} - C_{i,\mathrm{ocean}}\right)
\exp(-x/L_i).
\]

The crucial physical contrasts are:

\begin{itemize}
  \item \(T_{\mathrm{vent}} - T_{\mathrm{ocean}}\), mostly geophysics;
  \item \(pH_{\mathrm{vent}} - pH_{\mathrm{ocean}}\), mostly composition and
        water-rock chemistry;
  \item redox contrast between reduced vent fluid and ambient ocean.
\end{itemize}

\section{Titan Methane Pool Notes}

The Titan local environment should represent a methane pool, not water.
The pool should have:

\begin{itemize}
  \item methane solvent data;
  \item a time-dependent methane level;
  \item a spatial pool profile function stored in local environment data.
\end{itemize}

The current profile function lives in:

\begin{lstlisting}
src/environment/pool_spatial_profile.py
\end{lstlisting}

The plotting utility for this profile is:

\begin{lstlisting}
src/utilities/plot_titan_pool_profile.py
\end{lstlisting}

\section{Near-Term Implementation Order}

Recommended order:

\begin{enumerate}
  \item Finish the atmosphere solver structure: layers, stellar forcing, pressure
        and temperature profile.
  \item Produce a clear surface state from the atmosphere solver.
  \item Use \texttt{pH\_estimate.py} to estimate pH when water is present.
  \item Feed atmosphere-derived water state into
        \texttt{volcanic\_rocks/water\_level.py}.
  \item Add an aqueous chemistry backend driver, likely PHREEQC.
  \item Add hydrothermal vent inputs as ocean and vent reservoirs.
  \item Add a classifier that labels the final world type after solving.
\end{enumerate}

\end{document}
